\documentclass{beamer}
\usepackage{amsfonts,amsmath,oldgerm,mathrsfs,booktabs}
\usetheme{sintef}

\usepackage{xeCJK}
\usepackage{fontspec}
\setCJKmainfont{Noto Serif CJK SC}

\usefonttheme[onlymath]{serif}

% --- helpers ---
\newcommand{\takeaway}[1]{\begin{center}\large\bfseries #1\end{center}\vspace{-0.3em}}
\newcommand{\smallgap}{\vspace{0.5em}}

\titlebackground*{assets/background}

% -------- Title info --------
\title{Reproducing RoboDuet: A Cooperative Whole-body Controller for Legged Loco-Manipulation}
\subtitle{Progress Presentation (Supervisor Update)}
\course{Master's Degree in Artificial Intelligence}
\author{Yuxiao ZHAO}
\supervisor{Fei CHEN}
\IDnumber{1155246057}
\date{\today}

\begin{document}
\maketitle

% ==================================================
% SLIDE 1 — Project Scope
% ==================================================
\begin{frame}{Project Scope}
\takeaway{Reproduce and evaluate the RoboDuet controller in Isaac Gym simulation.}

\begin{columns}
\begin{column}{0.55\textwidth}
\textbf{Paper:} Pan et al., \emph{RoboDuet} (RAL 2025, arXiv:2403.17367)

\smallgap
\begin{itemize}
  \item Platform: \textbf{Isaac Gym}
  \item Robot: \textbf{Unitree B2 + Z1 arm}
  \item Key idea: two cooperative RL policies —
        a \emph{loco policy} for locomotion and
        an \emph{arm policy} for manipulation
  \item Policies interact: arm policy guides
        body posture via additional command outputs
\end{itemize}

\smallgap
\begin{block}{Out of Scope}
\small
\begin{itemize}
  \item No new algorithm design
  \item No sim-to-real transfer at this stage
\end{itemize}
\end{block}
\end{column}

\begin{column}{0.42\textwidth}
  % RoboDuet system overview diagram
  \includegraphics[width=\textwidth]{assets/overview}
  \small\centering
  \textit{RoboDuet two-policy architecture}
\end{column}
\end{columns}
\end{frame}

% ==================================================
% SLIDE 2 — Method Overview
% ==================================================
\begin{frame}{Method: Two-Stage Training in Isaac Gym}
\takeaway{Stage 1: robust locomotion. Stage 2: whole-body coordination.}

\begin{columns}
\begin{column}{0.56\textwidth}
\textbf{Stage 1 — Loco Policy}
\begin{itemize}
  \item PPO-based blind locomotion baseline
  \item Inputs: proprioception (joint pos/vel, IMU)
  \item Output: 12 leg joint position targets
  \item Trained until stable trotting gait emerges
\end{itemize}

\smallgap
\textbf{Stage 2 — Arm Policy (planned)}
\begin{itemize}
  \item Inputs: end-effector pose error + loco state
  \item Output: 6 arm joint targets +
        orientation command passed to loco policy
  \item Goal: 6-DoF end-effector pose tracking
        during locomotion
\end{itemize}
\end{column}

\begin{column}{0.42\textwidth}
\begin{block}{Cooperative Interaction}
\small
The arm policy outputs an orientation command
$a_t^{arm^G}$ that \emph{replaces} the body
orientation input of the loco policy, allowing
the body to adapt its posture for the arm.
\end{block}

\smallgap
\begin{block}{Training Environment}
\small
\begin{tabular}{@{}ll@{}}
Simulator  & Isaac Gym \\
Robot      & Unitree B2 + Z1 \\
Algorithm  & PPO \\
Randomisation & terrain, dynamics \\
\end{tabular}
\end{block}
\end{column}
\end{columns}
\end{frame}

% ==================================================
% SLIDE 3 — Current Progress
% ==================================================
\begin{frame}{Current Progress}
\takeaway{Stage 1 (locomotion baseline) reproduced in Isaac Gym.}

\begin{columns}
\begin{column}{0.55\textwidth}
\begin{block}{Completed}
\small
\begin{tabular}{@{}ll@{}}
\toprule
Item & Status \\
\midrule
Isaac Gym environment setup  & \textbf{Done} \\
Robot asset loaded (B2 + Z1) & \textbf{Done} \\
PPO locomotion baseline      & \textbf{Done} \\
Stable trotting gait         & \textbf{Done} \\
Arm policy training          & Ongoing \\
\bottomrule
\end{tabular}
\end{block}

\smallgap
\begin{block}{Locomotion Results}
\small
\begin{itemize}
  \item No fall over 3000-step rollout
  \item Stable forward locomotion achieved
  \item Reward converged after $\sim$X M steps
\end{itemize}
\end{block}
\end{column}

\begin{column}{0.42\textwidth}
  % Insert your training reward curve here
  \includegraphics[width=\textwidth]{assets/reward_curve}
  \small\centering
  \textit{PPO training reward curve (Stage 1)}
\end{column}
\end{columns}
\end{frame}

% ==================================================
% SLIDE 4 — Next Steps
% ==================================================
\begin{frame}{Next Steps}
\takeaway{Complete Stage 2 and demonstrate a pre-programmed grasp.}

\begin{enumerate}
  \item \textbf{Integrate arm DoFs} into policy action space
        (6 arm joints + orientation command output)
  \item \textbf{Train arm policy} (Stage 2) with
        6-DoF end-effector pose tracking objective
  \item \textbf{Evaluate whole-body coordination}:
        end-effector tracking accuracy, locomotion stability
  \item \textbf{Pre-programmed grasp demo}:
        fixed waypoint sequence for a simple pick task
        in Isaac Gym simulation
\end{enumerate}

\smallgap
\begin{alertblock}{Target Outcome}
\small
Simulation demo: quadruped walks to a target location
and executes a pre-programmed arm grasp sequence,
reproducing the core capability of RoboDuet.
\end{alertblock}
\end{frame}

% ==================================================
% THANK YOU
% ==================================================
\backmatter

\end{document}
